\section{Configuración y resolución de problemas al ejecutar MadGraph}
Esta sección se puede utilizar en caso de que tire problemas el \textit{launch} de un proceso.
\begin{itemize}
    \item[(1)] \textbf{Configuración de la interfaz LHAPDF}\\
    Dentro del entorno MG5
    \begin{lstlisting}
(MG5)> MG5_ROOT='/home/usuario/Descargas/madgraph/MG5_aMC_v2_9_27
(MG5)> LHAPDF_ROOT='$MG5_ROOT/HEPTools/lhapdf6_py3'

# Buscando la ubicacion del modulo LHAPDF
(MG5)> find "$LHAPDF_ROOT' -type d -name 'site-packages'

# Con la direccion se procede a acceder al archivo previamente creado
(MG5)> nano "$CONDA_PREFIX/etc/conda/activate.d/mg5_compiler.sh"

(mg5_compiler.sh)> export MG5_ROOT="/home/usuario/Descargas/madgraph/MG5_aMC_v2_9_27"
(mg5_compiler.sh)> export LHAPDF_ROOT="$MG5_ROOT/HEPTools/lhapdf6_py3"
(mg5_compiler.sh)> export LD_LIBRARY_PATH="$LHAPDF_ROOT/lib${LD_LIBRARY_PATH:+:$LD_LIBRARY_PATH}"
(mg5_compiler.sh)> export PYTHONPATH="$LHAPDF_ROOT/lib/python3.10/site-packages${PYTHONPATH:+:$PYTHONPATH}"

# Guardando y saliendo del entorno nano
(mg5_compiler.sh)> Ctrl + O
(mg5_compiler.sh)> Enter
(mg5_compiler.sh)> Ctrl + X

# Recargando el entorno MG5 y comprobando
(MG5)> conda deactivate
     > conda activate MG5
(MG5)> "$LD_LIBRARY_PATH"
(MG5)> echo "$PYTHONPATH"
(MG5)> python -c "import lhapdf; print(lhapdf.version())"
    \end{lstlisting}
    Una versión que verifique automaticamente la versión de Python, el cual busca el \textit{site-packages},
    \begin{lstlisting}
(MG5)> MG5_ROOT='/home/usuario/Descargas/madgraph/MG5_aMC_v2_9_27
(MG5)> LHAPDF_ROOT='$MG5_ROOT/HEPTools/lhapdf6_py3'

# Buscando la ubicacion del modulo LHAPDF
(MG5)> find "$LHAPDF_ROOT' -type d -name 'site-packages'

# Con la direccion se procede a acceder al archivo previamente creado
(MG5)> nano "$CONDA_PREFIX/etc/conda/activate.d/mg5_compiler.sh"

(mg5_compiler.sh)> export MG5_ROOT="/home/usuario/Descargas/madgraph/MG5_aMC_v2_9_27"
(mg5_compiler.sh)> export LHAPDF_ROOT="$MG5_ROOT/HEPTools/lhapdf6_py3"
(mg5_compiler.sh)> export LD_LIBRARY_PATH="$LHAPDF_ROOT/lib${LD_LIBRARY_PATH:+:$LD_LIBRARY_PATH}"
(mg5_compiler.sh)> LHAPDF_PYTHON_PATH="$(find "$LHAPDF_ROOT" - type d -name "site-packages" -print -quit)"
(mg5_compiler.sh)> if [ -n "$LHAPDF_PYTHON_PATH" ]; then
(mg5_compiler.sh)>    export PYTHONPATH="$LHAPDF_PYTHON_PATH${PYTHONPATH:+:$PYTHONPATH}"
(mg5_compiler.sh)> fi
(mg5_compiler.sh)> unset LHAPDF_PYTHON_PATH    
\end{lstlisting}
    \item[(2)] \textbf{Problema con \textit{Histogram with run\_id '0' xmldoc}}\\
    Este problema se da cuando pythia no puede leer los archivos xml, en caso de tener ese problema
    \begin{lstlisting}
(MG5)> unset PYTHIA8DATA
(MG5)> echo "$PYTHIA8DATA"
\end{lstlisting}
    \newpage 
    \begin{lstlisting}  
# Comprobando que se haya instalado Pythia en MadGraph
(MG5)> MG5_ROOT="/home/karna/madgraph/MG5_aMC_v2_9_27"
(MG5)> "$MG5_ROOT/HEPTools/pythia8/bin/pythia8-config" --version
<out> 8.317
(MG5)> "$MG5_ROOT/HEPTools/pythia8/bin/pythia8-config" --xmldoc
<out> /home/usuario/Descargas/madgraph/MG5_aMC_v2_9_27/HEPTools/pythia8/share/Pythia8/xmldoc

# Probando la ruta, indicando explicitamente los archivos XML
(MG5)> export PYTHIA8DATA="/home/usuario/Descargas/madgraph/MG5_aMC_v2_9_27/HEPTools/pythia8/share/Pythia8/xmldoc"
(MG5)> echo "$PYTHIA8DATA"
(MG5)> test -f "$PYTHIA8DATA/Index.xml" && echo "XML de Pythia encontrado"

# Verificando que la carpeta contiene la version correcta
(MG5)> grep -R "8.317" "$PYTHIA8DATA" | head

# Dejando la configuracion permanente en el entorno MG5, creando un script que se ejecute al final de los scripts de activacion
(MG5)> nano "$CONDA_PREFIX/etc/conda/activate.d/zz_mg5_pythia8.sh"

(zz_mg5_pythia8.sh)> export MG5_ROOT="/home/usuario/Descargas/madgraph/MG5_aMC_v2_9_27"
(zz_mg5_pythia8.sh)> export PYTHIA8DATA="$MG5_ROOT/HEPTools/pythia8/share/Pythia8/xmldoc"

# Guardando el script y retornando a la consola
(zz_mg5_pythia8.sh)> Ctrl + O
(zz_mg5_pythia8.sh)> Enter
(zz_mg5_pythia8.sh)> Ctrl + X

# Reactivando el entorno conda MG5 y verificando
(MG5)> conda deactivate
     > conda activate MG5 
(MG5)> echo "$PYTHIA8DATA"
<out> /home/usuario/Descargas/madgraph/MG5_aMC_v2_9_27/HEPTools/pythia8/share/Pythia8/xmldoc     
\end{lstlisting}
    En caso de que muestre esa salida ya está arreglada la configuración para archivos XML.
    \item[(3)] \textbf{Problema con la generación del archivo .root (Delphes)}\\
    En la terminal dentro de la carpeta de MadGraph (MG5\_aMC\_v2\_9\_27)
    \begin{lstlisting}
# Activando el entorno conda MG5 y verificando los archivos Delphes
     > conda activate MG5
(MG5)> find . -name "DelphesHepMC*" -ls
<out> ./Delphes/DelphesHepMC2
<out> ./Delphes/DelphesHepMC3
<out> ...
(MG5)> which root-config
<out> /home/usuario/anaconda3/envs/MG5/bin/root-config
(MG5)> root-config --version
<out> 6.40.02
(MG5)> root-config --prefix
<out> /home/usuario/anaconda3/envs/MG5
(MG5)> echo "$ROOTSYS"
<out> /home/usuario/anaconda3/envs/MG5

(MG5)> ln -s DelphesHepMC2 Delphes/DelphesHepMC
(MG5)> ls -l Delphes/DelphesHepMC*
<out> Delphes/DelphesHepMC -> DelphesHepMC2
<out> Delphes/DelphesHepMC2
<out> Delphes/DelphesHepMC3

# Verificando la resolucion
(MG5)> readlink -f Delphes/DelphesHepMC
<out> /home/usuario/Descargas/madgraph/MG5_aMC_v2_9_27/Delphes/DelphesHepMC2
\end{lstlisting}
\newpage
    \item[(4)] \textbf{Interfaz de MadGraph (launch)}\\
    Inicialmente al ejecutar el launch de un proceso se tendrá
\begin{lstlisting}
(MG5)> ./bin/mg5_aMC
MG5_aMC> launch "proceso"
 . . . . . . . . . . . . .
/===========================================================================\
| 1. Choose the shower/hadronization program     shower = OFF               |
| 2. Choose the detector simulation program    detector = OFF               |
| 3. Choose an analysis package (plot/convert) analysis = Not Avail.        |
| 4. Decay onshell particles                    madspin = OFF               |
| 5. Add weights to events for new hypp.       reweight = OFF               |
\===========================================================================/
# Para seleccionar un showering y un detector
MG5_aMC> shower=Pythia8
MG5_aMC> detector=Delphes
. . . . . . . . . . . . .
/===========================================================================\
| 1. Choose the shower/hadronization program     shower = Pythia8           |
| 2. Choose the detector simulation program    detector = Delphes           |
| 3. Choose an analysis package (plot/convert) analysis = Not Avail.        |
| 4. Decay onshell particles                    madspin = OFF               |
| 5. Add weights to events for new hypp.       reweight = OFF               |
\===========================================================================/
\end{lstlisting}
    \item[(5)] \textbf{Modelo a emplear en la simulación}\\
    Al ejecutar el entorno de MadGraph, tras la configuración previa, presenta la siguiente pantalla
\begin{lstlisting}
(MG5) usuario@Usuario:~/Descargas/madgraph/MG5_aMC_v2_9_27$ ./bin/mg5_aMC
************************************************************
*                     W E L C O M E to                     *
*              M A D G R A P H 5 _ a M C @ N L O           *
*                                                          *
*                 *                       *                *
*                   *        * *        *                  *
*                     * * * * 5 * * * *                    *
*                   *        * *        *                  *
*                 *                       *                *
*                                                          *
*         VERSION 2.9.27                2026-01-05         *
*    The MadGraph5_aMC@NLO Development Team - Find us at   *
*    https://server06.fynu.ucl.ac.be/projects/madgraph     *
*                            and                           *
*            http://amcatnlo.web.cern.ch/amcatnlo/         *
*               Type 'help' for in-line help.              *
*           Type 'tutorial' to learn how MG5 works         *
*    Type 'tutorial aMCatNLO' to learn how aMC@NLO works   *
*    Type 'tutorial MadLoop' to learn how MadLoop works    *
************************************************************
load MG5 configuration from input/mg5_configuration.txt 
set fastjet to fastjet-config
set lhapdf to /home/usuario/Descargas/madgraph/MG5_aMC_v2_9_27/HEPTools/lhapdf6_py3/bin/lhapdf-config
set lhapdf to /home/usuario/Descargas/madgraph/MG5_aMC_v2_9_27/HEPTools/lhapdf6_py3/bin/lhapdf-config
Using default eps viewer "evince". Set another one in ./input/mg5_configuration.txt
Using default web browser "firefox". Set another one in ./input/mg5_configuration.txt
Loading default model: sm
INFO: Restrict model sm with file models/sm/restrict_default.dat . 
INFO: Run "set stdout_level DEBUG" before import for more information. 
INFO: Change particles name to pass to MG5 convention 
Defined multiparticle p = g u c d s u~ c~ d~ s~
Defined multiparticle j = g u c d s u~ c~ d~ s~
Defined multiparticle l+ = e+ mu+
Defined multiparticle l- = e- mu-
Defined multiparticle vl = ve vm vt
Defined multiparticle vl~ = ve~ vm~ vt~
Defined multiparticle all = g u c d s u~ c~ d~ s~ a ve vm vt e- mu- ve~ vm~ vt~ e+ mu+ t b t~ b~ z w+ h w- ta- ta+
\end{lstlisting}
donde se puede observar el modelo cargado en la línea \textit{Loading default model: sm}. El modelo estándar sm presenta en el archivo \textit{param\_card.data} los siguientes valores
\begin{lstlisting}
  1 0.000000e+00 # d : 0.0 
  2 0.000000e+00 # u : 0.0 
  3 0.000000e+00 # s : 0.0 
  4 0.000000e+00 # c : 0.0 
  11 0.000000e+00 # e- : 0.0 
  12 0.000000e+00 # ve : 0.0 
  13 0.000000e+00 # mu- : 0.0 
  14 0.000000e+00 # vm : 0.0 
  16 0.000000e+00 # vt : 0.0 
  21 0.000000e+00 # g : 0.0 
  22 0.000000e+00 # a : 0.0 
  24 8.041900e+01 # w+ : cmath.sqrt(MZ\_\_exp\_\_2/2. + cmath.sqrt(MZ\_\_exp\_\_4/4. - (aEW*cmath.pi*MZ\_\_exp\_\_2)/(Gf*sqrt\_\_2)))
\end{lstlisting}
teniendo de esta forma una masa $m({\mu^{-}})=0.000000e+00$; por tanto, MadGraph trata al muón como no masivo en el elemento de matríz. Para cambiar esto, en el entorno de MadGraph se introduce el siguiente comando
\begin{lstlisting}
MG5_aMC> import model sm-full       # importando el modelo estandar completo
INFO: Change particles name to pass to MG5 convention 
Kept definitions of multiparticles p / j / l+ / l- / vl / vl~ unchanged
Defined multiparticle all = g u d s u~ d~ s~ a ve vm vt ve~ vm~ vt~ c t b c~ t~ b~ z w+ h w- e- mu- ta- e+ mu+ ta+
MG5_aMC> display particles mu-      # Verificando que mu- no tiene masa nula
Particle mu- has the following properties:
{
    'name': 'mu-',
    'antiname': 'mu+',
    'spin': 2,
    'color': 1,
    'charge': -1.00,
    'mass': 'mdl_MM',               # <- la masa de mu- es mdl\_MM
    'width': 'ZERO',
    'pdg_code': 13,
    'line': 'straight',
    'propagator': '',
    'is_part': True,
    'self_antipart': False,
    'type': '',
    'counterterm': {}
}
MG5_aMC> display parameters         # Verificando el valor de la masa de mu-
parameter type: ('external',)
        aEWM1 = 132.50698
        mdl_Gf = 1.16639e-05
        aS = 0.118
        mdl_lamWS = 0.2253
        mdl_AWS = 0.808
        mdl_rhoWS = 0.132
        mdl_etaWS = 0.341
        mdl_ymc = 1.27
        mdl_ymb = 4.2
        mdl_ymt = 164.5
        mdl_yme = 0.000511
        mdl_ymm = 0.10566
        mdl_ymtau = 1.777
        mdl_MZ = 91.188
        mdl_MC = 1.27
        mdl_MT = 172.0
        mdl_MB = 4.7
        mdl_MH = 125.0              # <- Valor de la masa de mu-
        mdl_Me = 0.000511
        mdl_MM = 0.10566
        mdl_MTA = 1.777
        mdl_WZ = 2.44140351
        mdl_WW = 2.04759951
        mdl_WT = 1.50833649
        mdl_WH = 0.00638233934
        mdl_WTau = 2.27e-12
\end{lstlisting}
Tras importar el modelo estándar completo, \textit{import model sm-full}, se puede generar el proceso
\begin{lstlisting}
# Primera version
MG5_aMC> import model sm-full
MG5_aMC> display particles mu-
MG5_aMC> display multiparticles
MG5_aMC> generate p p > mu+ mu-
MG5_aMC> display diagrams       # Comando para mostrar los diagramas de Feynman del proceso
MG5_aMC> output <archivo>

# Segunda version
MG5_aMC> import model sm-full
MG5_aMC> display particles mu-
MG5_aMC> display multiparticles
MG5_aMC> generate p p > mu+ mu- / h QED=2 QCD=0
MG5_aMC> display diagrams       # Comando para mostrar los diagramas de Feynman del proceso
MG5_aMC> output <archivo>
\end{lstlisting}
La segunda versión permite buscar el continuo Drell-Yan, además de evitar contribuciones de Higgs, dado que \textit{/ h} prohíbe diagramas internos con Higgs.
    \item[(6)] \textbf{Instalación de un paquete de análisis}\\
\end{itemize}